\scanpage{6}
\mspage{4}
\uline{Modèle d'une forme 1-dimensionnelle}

$L$, ensemble des ``\uline{lieux}''

$S$, \rule{3cm}{0.4pt} ``segments''

\[
S\longrightarrow\mathcal P(L),\qquad I\longmapsto\tilde I
\quad(\tilde I,\text{ ens des lieux \uline{sur} le segment }I)
\]
\[
S\longrightarrow\mathcal P_2(L)\qquad I\longmapsto\delta I
\quad(\delta I,\text{ ens des deux lieux \uline{extrémités} de }I)
\]
avec $\delta I\subset\tilde I$

$F_1)$ \unsure{Soient|Soit} $I\in S$, $\delta I=\{x,y\}$. \deleted{\unsure{Alors deux}}

[Soit $S_{I,x}$ l'ens des $J\in S$ tels que $\tilde J\subset\tilde I$, $\delta J\ni x$, et soit
$S_{I,x}\longrightarrow\tilde I\setminus\{x\}$ définie par
$J\longmapsto z$ tel que $\delta J=\{x,z\}$. Cette application est bijection, i.e.]
$\forall z\in\tilde I\setminus x$, $\exists!\ J\in S$, \deleted{\unsure{un}} soit $I_{x,z}$ tel que
\[
\unsure{\tilde J_{x,z}|\tilde I_{x,z}}\subset\tilde I,\qquad
\unsure{\delta J_{x,z}|\delta I_{x,z}}=\{x,z\}
\]
\margin{axiome de décomposition}
\margin{NB $I_{x,y}=I_{y,x}=I$}

\deleted{De plus, on a si $z\in\tilde I^\circ$ (i.e. $z\in\tilde I\setminus\{x,y\}$),}
\inserted{$F_2)$ Soit $I\in S$, alors $\tilde I^\circ\neq\varnothing$, et si $z\in\tilde I^\circ$,}
\deleted{on a}
\[
\tilde I_{x,z}\cap\tilde I_{y,z}=\{z\}.
\]
\margin{(\unsure{7.6.|6.6.}) \unsure{Faudra|Faut} \unsure{indiquer} qu'on \unsure{fait sauter|fait passer} la condition de divisibilité $\tilde I^\circ\neq\varnothing$.}

Si $J\in S$ est tel que $\tilde J\subset\tilde I$, $z\in\delta J$,
on a $\tilde J\subset\tilde I_{x,z}$ \uline{ou} $J\subset\tilde I_{y,z}$.

\uline{Cor. 1} \deleted{$\unsure{\forall}$} $\forall I\in S$, $\tilde I^\circ$ est infini [du fait de la condition de divisibilité)

\uline{Cor. 2} Dans $F_2$, les relations $\tilde J\subset\tilde I_{x,z}$ et $\tilde J\subset\tilde I_{y,z}$ sont exclusives l'une de l'autre.

\uline{Cor 3} Soit $I\in S$, \deleted{$\unsure{\exists}$} $S_I=\{J\in S\mid\tilde J\subset\tilde I\}$, et
\[
S_I\unsure{-|\longrightarrow}\mathcal P_2(\tilde I)\qquad J\longmapsto\delta J.
\]
Alors \deleted{\unsure{cette}}

a) Cette application est injective.

b) Son image contient les $\varepsilon\in\mathcal P_2(\tilde I)$ tels que $\varepsilon\cap\delta I\neq\varnothing$.
\mspage{5}

\rotatebox{90}{\reflectbox{\unsure{EXTRA STRONG}}}
\rotatebox{90}{\reflectbox{\unsure{EXTRA STRONG}}}
